\documentclass{ctexart}
\begin{document}
\begin{abstract}
本文建立几何模型并报告结果、约束与检验。
\end{abstract}
\keywords{几何，误差，稳健性}
\section{问题重述}
题目要求相邻节点距离变化率不得超过 $0.07\%$。
\section{问题分析}
采用曲面几何与连续参数优化分析节点调整。
\section{模型假设}
假设测量数据有效，并在结果中检查假设影响。
\section{符号说明}
$d_i$ 表示第 $i$ 个节点的相对距离变化率。
\section{分问题模型建立与求解}
模型一按最小化最大偏差求解。计算显示仅有 $54.94\%$ 的节点满足
$0.07\%$ 阈值。模型二没有求解，因此只列为备选思路。
接收率使用经验修正系数 2.899 进行计算。
\section{结果分析}
结果表明目标值相对基准有所变化，但硬约束仍需人工裁决和修复。
每个参数取值均得到一个稳定的数值，因此原稿据此声称模型稳定。
\section{模型检验}
通过残差、基准对照和 RMSE 检查数值结果。
\section{灵敏度与稳健性分析}
对顶点位置作参数扰动并记录响应曲线。
\section{模型评价与改进}
模型可用于几何候选筛选，后续需要补充结构约束和物理接收链。
\begin{thebibliography}{1}
\bibitem{sample} 示例文献。
\end{thebibliography}
\appendix
\section{附录}
程序入口和参数记录见随文文件。
\end{document}
